\subsection{Functions with values in a space of measures}

Let X be a locally compact space, \(\mathcal{M}_{+}(X)\) the convex cone of positive measures on X. Throughout the rest of this chapter, \(\mathcal{M}_{+}(X)\) will be equipped with the topology induced by the vague topology on \(\mathcal{M}(X)\) (Ch. III, §1, No.~9); thus, to say that a mapping \(\Lambda: t \mapsto \lambda_{t}\) of the locally compact space T into \(\mathcal{M}_{+}(X)\) is continuous means that, for every function \(f \in \mathcal{K}(X)\) , the numerical function \(t \mapsto \lambda_{t}(f)\) is continuous. In this case we shall also say that \(\Lambda\) is vaguely continuous on T. To say that a mapping \(\Lambda: t \mapsto \lambda_{t}\) is \(\mu\) -measurable means that the set of compact subsets K of T, such that the restriction of \(\Lambda\) to K is vaguely continuous, is \(\mu\) -dense (Ch. IV, §5, No.~10, Prop. 15). We shall then say that \(\Lambda\) is vaguely \(\mu\) -measurable.

Let \(\Lambda : t \mapsto \lambda_t\) be a mapping of T into \(\mathcal{M}_+(\mathrm{X})\) ; we shall say that \(\Lambda\) is scalarly essentially integrable for the measure \(\mu\) if, for every function \(f \in \mathcal{K}(\mathrm{X})\) , the function \(t \mapsto \lambda_t(f)\) is essentially \(\mu\) -integrable. If one sets \(\nu(f) = \int \lambda_t(f) d\mu(t)\) , it is clear that \(\nu\) is a positive linear form on \(\mathcal{K}(\mathrm{X})\) , hence is a measure on X (Ch. III, §1, No.~5, Th. 1). We will say that \(\nu\) is the integral of the function \(\Lambda\) with values in \(\mathcal{M}_+(\mathrm{X})\) , and we will write \(\nu = \int \lambda_t d\mu(t)\) .

The preceding definition is a special case of the concept of weak integral, which will be treated in a general manner in Ch. VI.

If \(f\) denotes an element of \(\mathcal{K}(\mathrm{X})\) , the integral \(\int \lambda_t(f) d\mu(t)\) will also, by an abuse of notation, be denoted \(\int d\mu(t) \int f(x) d\lambda_t(x)\) ; the definition of the integral \(\nu = \int \lambda_t d\mu(t)\) may then be written

\[
\int f (x) d \nu (x) = \int d \mu (t) \int f (x) d \lambda_ {t} (x).\tag{1}
\]

We shall make analogous abuses of notation in the sequel, for upper integrals, essential upper integrals, and integrals of functions with values in a Banach space.

Examples. --- 1) Suppose that T is a discrete space, and that \(\mu\) is the measure on T defined by placing a mass \(+1\) at each point of T (Ch. III, §1, No.~3). Let \(h\) be a function \(\geqslant 0\) defined on T; since the function \(h\) is lower semi-continuous (even continuous) on T, \(\mu^{*}(h) = \mu^{\bullet}(h) = \sum_{t\in \mathrm{T}}h(t)\) (Ch. IV, §1, No.~1, Example).

For the measure \(\mu\) , the notions of integrable function and essentially integrable function are therefore identical. This being so, to say that a mapping \(t \mapsto \lambda_t\) of T into \(\mathcal{M}_+(X)\) is scalarly essentially \(\mu\) -integrable amounts to saying that the family \((\lambda_t)_{t \in \mathrm{T}}\) is summable (§2, No.~1), and one then has \(\int \lambda_t d\mu(t) = \sum_{t \in \mathrm{T}} \lambda_t\) . Note that the mapping \(t \mapsto \lambda_t\) is vaguely continuous.

\begin{enumerate}
\def\labelenumi{\arabic{enumi})}
\setcounter{enumi}{1}
\tightlist
\item
  The mapping \(t \mapsto \varepsilon_t\) of \(T\) into \(\mathcal{M}_+(T)\) is vaguely continuous, scalarly essentially \(\mu\) -integrable for every positive measure \(\mu\) on \(T\) , and one has \(\int \varepsilon_t d\mu(t) = \mu\) .
\end{enumerate}

PROPOSITION 1. --- Suppose that \(\mu\) is the supremum of an increasing directed family \((\mu_i)_{i \in I}\) of positive measures on T; in order that \(\Lambda : t \mapsto \lambda_t\) be scalarly essentially \(\mu\) -integrable, it is necessary and sufficient that \(\Lambda\) be scalarly essentially \(\mu_i\) -integrable for all \(i \in I\) and that the family \(\left( \int \lambda_t d\mu_i(t) \right)_{i \in I}\) be bounded above in \(\mathcal{M}(X)\) . One then has,

\[
\int \lambda_ {t} d \mu (t) = \sup _ {i \in \mathrm{I}} \int \lambda_ {t} d \mu_ {i} (t).\tag{2}
\]

For, verifying that \(\Lambda\) is scalarly essentially integrable for a positive measure \(\eta\) on T comes down to verifying that \(t \mapsto \lambda_{t}(g)\) is \(\eta\) -measurable and admits a finite essential upper integral, with respect to \(\eta\) , for every function \(g \in \mathcal{K}_{+}(X)\) . The proposition therefore follows at once from Prop. 11 of §1, No.~4 and its Corollary 2.

COROLLARY. --- Suppose that \(\mu\) is the sum of a summable family \((\mu_{\alpha})_{\alpha \in \mathbf{A}}\) of positive measures on \(\mathbf{T}\) ; in order that \(\Lambda : t \mapsto \lambda_t\) be scalarly essentially \(\mu\) -integrable, it is necessary and sufficient that \(\Lambda\) be scalarly essentially \(\mu_{\alpha}\) -integrable for every \(\alpha \in \mathbf{A}\) and that the family of measures \(\int \lambda_t d\mu_{\alpha}(t)\) be summable. One then has

\[
\int \lambda_ {t} d \mu (t) = \sum_ {\alpha \in A} \int \lambda_ {t} d \mu_ {\alpha} (t).\tag{3}
\]

It follows immediately that every scalarly essentially \(\mu\) -integrable mapping is also scalarly essentially \(\mu'\) -integrable for every measure \(\mu' \leqslant \mu\) .

In this section we shall limit ourselves to the study of scalarly essentially integrable mappings of T into \(\mathcal{M}_{+}(X)\) that have the property contemplated in the following definition:

DEFINITION 1. --- Let X be a locally compact space, \(\Lambda : t \mapsto \lambda_t\) a scalarly essentially \(\mu\) -integrable mapping of T into \(\mathcal{M}_{+}(X)\) , and \(\nu\) the integral of \(\Lambda\) .

We say that \(\Lambda\) is \(\mu\) -pre-adequate if, for every lower semi-continuous function \(f \geqslant 0\) defined on X, the function \(t \mapsto \int^{\bullet} f d\lambda_t\) is \(\mu\) -measurable on T and

\[
\int^ {\bullet} f (x) d \nu (x) = \int^ {\bullet} d \mu (t) \int^ {\bullet} f (x) d \lambda_ {t} (x).\tag{4}
\]

We say that \(\Lambda\) is \(\mu\) -adequate (*) if \(\Lambda\) is \(\mu'\) -pre-adequate for every positive measure \(\mu' \leqslant \mu\) .

It can be shown that if \(\Lambda\) is \(\mu\) -pre-adequate and if the measure \(\nu = \int \lambda_{t} d\mu(t)\) is moderated---in particular if X is countable at infinity---then \(\Lambda\) is \(\mu\) -adequate (Exer. 7); however, it is not known if these concepts are in general equivalent. In the statements in Nos. 2 and 3 below, the assertions preceded by an a) or a b) extend at once to pre-adequate mappings, whereas those preceded by a c) are valid only for adequate mappings.

The following proposition often permits one to verify that a given mapping is \(\mu\) -adequate.

PROPOSITION 2. --- Let \(\Lambda : t \mapsto \lambda_t\) be a scalarly essentially \(\mu\) -integrable mapping of T into \(\mathcal{M}_{+}(X)\) , and let \(\nu = \int \lambda_t d\mu(t)\) .

\begin{enumerate}
\def\labelenumi{\alph{enumi})}
\tightlist
\item
  If \(\Lambda\) is vaguely continuous, then the mapping \(t \mapsto \lambda_t^\bullet(f)\) is lower semi-continuous for every lower semi-continuous function \(f \geqslant 0\) defined on X, \(\Lambda\) is \(\mu\) -adequate, and we have the relation
\end{enumerate}

\[
\int^ {*} f (x) d \nu (x) = \int^ {*} d \mu (t) \int^ {*} f (x) d \lambda_ {t} (x).\tag{5}
\]

\begin{enumerate}
\def\labelenumi{\alph{enumi})}
\setcounter{enumi}{1}
\item
  If \(\Lambda\) is vaguely \(\mu\) -measurable, then \(\Lambda\) is \(\mu\) -adequate.
\item
  If the topology of X admits a countable base, then \(\Lambda\) is vaguely \(\mu\) -measurable (hence also \(\mu\) -adequate).
\end{enumerate}

Let \(f\) be a lower semi-continuous function \(\geqslant 0\) defined on \(X\) . Let \(F\) be the set, directed for the relation \(\leqslant\) , of functions \(g \in \mathcal{H}(X)\) such that \(0 \leqslant g \leqslant f\) . For \(g \in F\) , denote by \(h_g\) the function defined on \(T\) by \(h_g(t) = \lambda_t(g)\) . Similarly, set

\[
h _ {f} (t) = \lambda_ {t} ^ {*} (f) = \lambda_ {t} ^ {\bullet} (f) = \sup _ {g \in \mathrm{F}} h _ {g} (t)
\]

(§1, No.~1, Prop. 4). Let us make the following hypothesis, weaker than that of a): assume only that the restriction of \(\Lambda\) to S is vaguely continuous, where S is a closed subset of T that contains the support of \(\mu\) . For \(g \in \mathbf{F}\) , denote by \(\overline{h}_g\) the numerical function that coincides with \(h_g\) on S and has the value \(+\infty\) on \(\mathbb{C}S\) . Set \(\overline{h}_f = \sup_{g \in \mathbb{F}} \overline{h}_g\) ; then \(\overline{h}_f = h_f\) on S. For every \(g \in \mathbf{F}\) , the function \(\overline{h}_g\) is lower semi-continuous; \(\overline{h}_f\) is therefore lower semi-continuous and, since the family \((\overline{h}_g)_{g \in \mathbb{F}}\) is directed,

\[
\mu^ {*} (\overline {{h}} _ {f}) = \sup _ {g \in \mathrm{F}} \mu^ {*} (\overline {{h}} _ {g}) = \sup _ {g \in \mathrm{F}} \mu^ {*} (h _ {g}) = \sup _ {g \in \mathrm{F}} \nu (g) = \nu^ {*} (f)
\]

(Ch. IV, §1, No.~1, Th. 1 and §2, No.~3, Prop. 6). Since \(h_f = \overline{h}_f\) on S, hence almost everywhere, this may also be written \(\mu^*(h_f) = \nu^*(f)\) , an equality identical to (5). Similarly, \(f\) and \(\overline{h}_f\) being lower semi-continuous, the preceding relations yield the equality \(\mu^\bullet(\overline{h}_f) = \nu^\bullet(f)\) (§1, No.~1, Prop. 4); since \(\overline{h}_f = h_f\) on S, it follows that \(\mu^\bullet(h_f) = \nu^\bullet(f)\) (§1, No.~1, Prop. 1), an equality identical to (4). The mapping \(\Lambda\) is therefore \(\mu\) -pre-adequate; but one could have replaced everywhere in this argument \(\mu\) by \(\mu' \leqslant \mu\) , and \(\nu\) by \(\nu' = \int \lambda_t d\mu'(t)\) , because \(\Lambda\) is also scalarly essentially \(\mu'\) -integrable and S contains the support of \(\mu'\) . It follows from this that \(\Lambda\) is \(\mu\) -adequate.

Suppose \(\Lambda\) is vaguely continuous; we may take \(S = T\) ; then \(h_f = \overline{h}_f\) is lower semi-continuous, which completes the proof of part a) of the statement.

Assume that \(\Lambda\) is vaguely \(\mu\) -measurable and let us prove b). The set \(\mathfrak{K}\) of compact subsets K of T such that the restriction of \(\Lambda\) to K is continuous being \(\mu\) -dense (Ch. IV, §5, No.~10, Prop. 15), there exists a summable family \((\mu_{\alpha})_{\alpha \in \mathrm{A}}\) of measures on T such that \(\mu = \sum_{\alpha \in \mathrm{A}} \mu_{\alpha}\) and the support of each of the measures \(\mu_{\alpha}\) belongs to \(\mathfrak{K}\) (§2, No.~3, Prop. 4). For every \(\alpha \in \mathrm{A}\) , the mapping \(\Lambda\) is scalarly essentially \(\mu_{\alpha}\) -integrable, and we set \(\nu_{\alpha} = \int \lambda_t d\mu_{\alpha}(t)\) ; the family \((\nu_{\alpha})\) is summable, and its sum is equal to \(\nu\) (Cor. of Prop. 1). If \(f\) is a positive lower semi-continuous function defined on X, the first part of the proof, applied to the measures \(\mu_{\alpha}\) and the closed sets \(S_{\alpha}\) , shows that:

\(1^{\circ} h_{f}\) is \(\mu_{\alpha}\) -measurable for every \(\alpha \in A\) , hence is \(\mu\) -measurable (§2, No.~2, Prop. 2), and

\[
2 ^ {\circ} \int^ {\bullet} f (x) d \nu_ {\alpha} (x) = \int^ {\bullet} d \mu_ {\alpha} (t) \int^ {\bullet} f (x) d \lambda_ {t} (x).
\]

The formula (4) follows on summing over \(\alpha\) ( \(\S 2\) , No.~2, Prop. 1). Applying the preceding argument to an arbitrary measure \(\mu' \leqslant \mu\) (which is legitimate, since \(\Lambda\) is scalarly essentially \(\mu'\) -integrable and vaguely \(\mu'\) -measurable, cf.~\(\S 2\) , No.~2, Prop. 2), we conclude that \(\Lambda\) is \(\mu'\) -pre-adequate, and b) is proved.

Finally, assuming that the topology of X admits a countable base, let us show that every scalarly essentially \(\mu\) -integrable mapping \(\Lambda: t \mapsto \lambda_{t}\) of T into \(\mathcal{M}_{+}(X)\) is vaguely \(\mu\) -measurable. This will result from the following lemma:

Lemma. --- Let X be a locally compact space having a countable base. Then, there exists in \(\mathcal{K}(\mathrm{X})\) a countable subset S having the following property: for every function \(f \in \mathcal{K}(\mathrm{X})\) , there exist a sequence \((f_n)\) of elements of S and a positive function \(\varphi \in S\) such that, for every number \(\varepsilon > 0\) , \(|f_n - f| \leqslant \varepsilon \varphi\) provided n is sufficiently large.

Let \(X'\) be the Alexandroff compactification of \(X\) , which is a metrizable compact space (GT, IX, §2, No.~9, Prop. 16 and Cor.); we identify \(\mathcal{K}(X)\) with a subset of \(\mathcal{C}(X')\) . Let \(S'\) be a countable dense subset of the Banach space \(\mathcal{C}(X')\) (GT, X, §3, No.~3, Th. 1); we can suppose that \(S'\) contains the constant function \(n\) for every \(n \in N\) . Let \((U_n)\) be a sequence of relatively compact open sets in \(X\) , with union \(X\) , such that \(\overline{U}_n \subset U_{n+1}\) for all \(n\) (GT, I, §9, No.~9, Prop. 15), and let \(\varphi_n\) be a function in \(\mathcal{K}_+(X)\) equal to 1 on \(\overline{U}_n\) . We denote by \(S\) the countable set of elements of \(\mathcal{K}(X)\) of the form \(\varphi_n g\) ( \(n \in N, g \in S'\) ). If \(f \in \mathcal{K}(X)\) , let \((g_n)\) be a sequence of elements of \(S'\) that converges uniformly to \(f\) , and let \(k\) be an integer such that the support of \(f\) is contained in \(U_k\) . Finally, let \(m\) be an integer that is an upper bound for the norms of the functions \(g_n\) . The functions \(f_n = \varphi_k g_n\) belong to \(S\) and satisfy the statement, with \(\varphi = m\varphi_k\) .

The lemma having been established, and the mapping \(t \mapsto \lambda_t(g)\) being essentially \(\mu\) -integrable for every \(g \in S\) , the mapping \(t \mapsto (\lambda_t(g))_{g \in S}\) of \(T\) into \(\mathbf{R}^S\) is \(\mu\) -measurable (Ch. IV, §5, No.~3, Th. 1). The set \(\mathfrak{K}\) , of compact subsets \(K\) of \(T\) such that the restriction of this mapping to \(K\) is continuous, is therefore \(\mu\) -dense, and it will suffice to show that the restriction of \(\Lambda\) to every \(K \in \mathfrak{K}\) is continuous. Now, let \(f\) be any element of \(\mathcal{H}(X)\) , \(f_n\) and \(\varphi\) elements of \(S\) satisfying the statement of the Lemma; the function \(t \mapsto \lambda_t(f)\) is then the uniform limit on \(K\) of the continuous functions \(t \mapsto \lambda_t(f_n)\) ; it is therefore continuous on \(K\) , and the proposition is proved.
% semantic-fragments: legacy-input-seam
\relax{}
